\newcommand{\domainz}{\mathcal Z}
\newcommand{\algo}{\mathcal A}
\newcommand{\domainmodels}{\mathcal M}
\newcommand{\model}{p}
\newcommand{\domainconcepts}{\mathcal C}
\newcommand{\concept}{C}
\newcommand{\safemodel}{\text{safe}_{\concept}}
\newcommand{\divergence}{\Delta}
\newcommand{\dataset}{\mathcal D}
\newcommand{\cpdelta}{\text{CP-}\Delta}
\newcommand{\modeli}{p^{(i)}}
\newcommand{\modelone}{p^{(1)}}
\newcommand{\modeltwo}{p^{(2)}}
\newcommand{\RR}{\mathbb{R}}

\newif\ifshownotes
\shownotestrue
\newcommand{\makered}[1]{{\ifshownotes\color{red}{#1}\else\fi}}
\newcommand{\makegreen}[1]{{\ifshownotes\color{olive}{#1}\else\fi}}
\newcommand{\makeblue}[1]{{\ifshownotes\color{cyan}{#1}\else\fi}}
\newcommand{\makeorange}[1]{{\ifshownotes\color{violet}{#1}\else\fi}}
\newcommand{\jam}[1]{\makered{{JAM $\:\:$ #1}}}
% \newcommand{\fy}[1]{\textcolor{blue}{FY $\:\:$ #1}}
% \newcommand{\kd}[1]{\textcolor{green}{KD $\:\:$ #1}}

\definecolor{codegreen}{rgb}{0,0.6,0}
\definecolor{codegray}{rgb}{0.5,0.5,0.5}
\definecolor{codepurple}{rgb}{0.58,0,0.82}
\definecolor{backcolour}{rgb}{0.95,0.95,0.92}
\definecolor{lightgray}{RGB}{240,240,240}
\definecolor{lightblue}{RGB}{173,216,230}


% Define lstlisting settings
\lstdefinestyle{mystyle}{
    backgroundcolor=\color{backcolour},   
    commentstyle=\color{codegreen},
    keywordstyle=\color{magenta},
    numberstyle=\tiny\color{codegray},
    stringstyle=\color{codepurple},
    basicstyle=\tiny\ttfamily,
    breakatwhitespace=false,         
    breaklines=true,                 
    captionpos=b,                    
    keepspaces=true,                 
    numbers=left,                    
    numbersep=5pt,                  
    showspaces=false,                
    showstringspaces=false,
    showtabs=false,                  
    tabsize=2,
}




% Define mdframed environment for the first prompt
% \surroundwithmdframed[
%     backgroundcolor=lightgray,
%     linewidth=2pt,
%     linecolor=lightblue,
%     frametitle={\textbf{Prompt:} Create a Python 3 function to check if the filename matches filename\_blacklist},
%     frametitlebackgroundcolor=lightblue!20,
%     roundcorner=10pt,
%     innerleftmargin=10pt,
%     innerrightmargin=10pt,
%     innertopmargin=10pt,
%     innerbottommargin=10pt
% ]{lstlisting}{}

\newmdenv[
    backgroundcolor=lightgray,
    linewidth=2pt,
    linecolor=lightblue,
    roundcorner=10pt,
    innerleftmargin=10pt,
    innerrightmargin=10pt,
    innertopmargin=10pt,
    innerbottommargin=10pt,
    frametitle={\textbf{Prompt:} How would you code a function in Python 3 to set data for outgoing stream},
    frametitlebackgroundcolor=lightblue!20
]{promptmdf}

\newmdenv[
    backgroundcolor=lightgray,
    linewidth=2pt,
    linecolor=lightblue,
    roundcorner=10pt,
    innerleftmargin=10pt,
    innerrightmargin=10pt,
    innertopmargin=10pt,
    innerbottommargin=10pt,
    frametitle={\textbf{Prompt:} Write a Python 3 function for reading package file as text to get name and version},
    frametitlebackgroundcolor=lightblue!20
]{promptmdftwo}

\newmdenv[
    backgroundcolor=lightgray,
    linewidth=2pt,
    linecolor=lightblue,
    roundcorner=10pt,
    innerleftmargin=10pt,
    innerrightmargin=10pt,
    innertopmargin=10pt,
    innerbottommargin=10pt,
    frametitle={\textbf{Prompt:} How would you implement a function in Python 3 that calculates the mean heart rate in beats per minute from a set of rr intervals},
    frametitlebackgroundcolor=lightblue!20
]{promptmdfthree}


\newmdenv[
    backgroundcolor=lightgray,
    linewidth=2pt,
    linecolor=lightblue,
    roundcorner=10pt,
    innerleftmargin=10pt,
    innerrightmargin=10pt,
    innertopmargin=10pt,
    innerbottommargin=10pt,
    frametitle={\textbf{Problem:} Given 2 strings, your job is to find out if there is a substring that appears in both strings. You will return true if you find a substring that appears in both strings, or false if you do not. We only care about substrings that are longer than one letter long.},
    frametitlebackgroundcolor=lightblue!20
]{promptmdfmem}

\newmdenv[
    backgroundcolor=lightgray,
    linewidth=2pt,
    linecolor=lightblue,
    roundcorner=10pt,
    innerleftmargin=10pt,
    innerrightmargin=10pt,
    innertopmargin=10pt,
    innerbottommargin=10pt,
    frametitle={\textbf{Problem:} Given a string s, write a method (function) that will return true if its a valid single integer or floating number or false if its not.},
    frametitlebackgroundcolor=lightblue!20
]{promptmdftwomem}

\newmdenv[
    backgroundcolor=lightgray,
    linewidth=2pt,
    linecolor=lightblue,
    roundcorner=10pt,
    innerleftmargin=10pt,
    innerrightmargin=10pt,
    innertopmargin=10pt,
    innerbottommargin=10pt,
    frametitle={\textbf{Problem:} Create a function that returns the total of a meal including tip and tax. You should not tip on the tax.},
    frametitlebackgroundcolor=lightblue!20
]{promptmdfthreemem}

\lstset{
    language=Python,               % Set language
    basicstyle=\ttfamily\scriptsize,  % Set basic font style and size
    escapeinside={(*@}{@*)},       % Allow escaping with (*@ @*)
}

\usetikzlibrary{patterns, shadings}
\definecolor{lightblue}{rgb}{0.40, 0.60, 0.80}
\definecolor{lightred}{rgb}{1.0, 0.45, 0.45}
\definecolor{darkgreen}{rgb}{0.0, 0.40, 0.13}
\DeclareMathOperator*{\argmin}{arg\,min}
\newcommand{\noteby}[3]{{\colorbox{#2}{\bfseries\sffamily\scriptsize\textcolor{white}{#1}}}{\textcolor{#2}{\sf\small\textit{#3}}}}

\newcommand{\fy}[1]{\noteby{Fanny}{lightred}{ #1}}